\chapter{Order-one Krajewski square: факторизация, \texorpdfstring{\(J\)}{J} и CP-фаза}

Внешний red-team справедливо потребовал вычислить
\begin{equation}
 [[D_F,a],JbJ^{-1}]
\end{equation}
до объявления connector graph допустимой finite geometry. Выполним этот
тест на минимальном summand-label menu, не предполагая заранее исчезновение
CP-фазы.

\section{Точный edge criterion}

Обозначим неприводимый бимодуль через
\(\mathcal H_{ij}\), где \(i\) задаёт left action, а \(j\) --- right action.
Для блока \(D_{ij,kl}\) double commutator содержит множитель
\begin{equation}
 (a_i-a_k)D_{ij,kl}(b_j-b_l).
\end{equation}
Поэтому order-one condition для всех \(a,b\) разрешает ненулевой edge только
если
\begin{equation}
 i=k
 \qquad\text{или}\qquad
 j=l.
\end{equation}
Иными словами, один Dirac edge может изменить left label или right label,
но не оба одновременно.

Для четырёх coarse nodes
\begin{equation}
 A=(o,o),\quad B=(o,h),\quad C=(h,h),\quad D=(h,o)
\end{equation}
разрешены ровно edges
\begin{equation}
 AB,\quad BC,\quad CD,\quad DA,
\end{equation}
а diagonals \(AC\) и \(BD\) запрещены.

\begin{theorem}[Order-one rectangle theorem]
Треугольный observed--messenger--hidden contour не является допустимым
Krajewski cycle, если хотя бы один его edge одновременно меняет left и right
algebra labels. Минимальный connector cycle, проходящий order-one condition,
есть четырёхрёберный rectangle.
\end{theorem}

Таким образом, EFT-chain предыдущей главы не должна буквально отождествляться
с тремя finite-Dirac edges. Её spectral lift обязан разрешиться в alternating
horizontal/vertical square. Это подтверждает факторизацию gauge intertwiner
и family multiplicity, отмеченную red-team.

\section{\texorpdfstring{\(J\)}{J}-совместимый квадрат}

Реальная структура действует как
\begin{equation}
 J:A\mapsto A,\qquad B\leftrightarrow D,\qquad C\mapsto C.
\end{equation}
После допустимых rephasings минимальный self-adjoint \(J\)-совместимый
оператор можно записать через два модуля \(a,b>0\) и одну loop phase
\(\phi\). Прямой matrix calculation даёт
\begin{align}
 \Tr D_F^2 &=4a^2+4b^2,\\
 \Tr D_F^4 &=8a^4+8b^4+16a^2b^2\cos^2\phi,\\
 \det D_F &=4a^2b^2\sin^2\phi.
\end{align}

\begin{proposition}[Фаза не уничтожается реальной структурой]
\(J\)-идентификация противоположных edges оставляет одну
rephasing-invariant loop phase. Для finite potential
\begin{equation}
 V_F=-\mu^2\Tr D_F^2+\lambda\Tr D_F^4,
 \qquad \lambda>0,
\end{equation}
phase minima равны
\begin{equation}
 \phi=\pm\frac{\pi}{2}.
\end{equation}
Следовательно, order-one factorization сама по себе не убивает maximal
complex phase.
\end{proposition}

Однако вещественное spectral action зависит от \(\cos^2\phi\) и не различает
две ориентации. Оно допускает spontaneous CP branch, но не предсказывает знак
orientation.

\section{Где obstruction действительно остаётся}

Family space является multiplicity factor, коммутирующим с gauge algebra.
Поэтому разрешённый edge имеет форму
\begin{equation}
 D_e=I_e^{\rm gauge}\otimes Y_e^{\rm family}.
\end{equation}
Order-one condition не ограничивает \(Y_e\) до rank-one projector или
incidence operator. Если хотя бы один \(Y_e\) произволен, а остальные path
factors обратимы, effective path product пробегает всю
\(M_3(\mathbb C)\). Машинный basis test даёт span dimension \(9\).

\begin{nogocriterion}[Family-factor freedom]
Order-one pass и surviving loop phase не закрывают flavour sector. Connector
считается predictive только если тот же parent action ограничивает family
edge factors заранее объявленной алгеброй, например выбранными
\(P_-\) и incidence directions, без произвольного \(M_3\)-входа.
\end{nogocriterion}

Итог уточняет внешний аудит: доказанная ошибка исходного sketch --- не сама
факторизация и не автоматическая гибель CP, а использование contour до
явного order-one lift. Rectangle lift существует и сохраняет фазу; открытым
остаётся selector семейных matrices и ориентационный знак. Машинная проверка
сохранена в
\path{s2t_v4_order_one_krajewski_square_gate_results.json}.